\documentclass[11pt]{article}
\usepackage[margin=1in]{geometry}
\usepackage{amsmath,amssymb,amsthm,mathtools,booktabs}
\usepackage[hidelinks]{hyperref}

\newtheorem{theorem}{Theorem}
\newtheorem{lemma}{Lemma}
\newtheorem{proposition}{Proposition}
\newtheorem{corollary}{Corollary}
\newtheorem{remark}{Remark}

\newcommand{\ip}[2]{\left\langle #1,#2\right\rangle}
\newcommand{\norm}[1]{\left\lVert #1\right\rVert}
\newcommand{\R}{\mathbb R}
\newcommand{\D}{\mathrm D}
\newcommand{\dd}{\mathrm d}

\title{Bregman-Balanced Integrators, Discrete Work, and Exact Geodesic Balance Defects}
\author{}
\date{}

\begin{document}
\maketitle

\section{Setup}

Let $\phi$ be a strictly convex $C^3$ barrier on a convex domain, and write
\[
    g(x):=\nabla\phi(x),
    \qquad
    H(x):=\nabla^2\phi(x).
\]
The Hessian metric defines
\[
    \ip{p}{q}_x:=p^\top H(x)q,
    \qquad
    \norm{p}_x^2:=\ip{p}{p}_x,
    \qquad
    \norm{w}_{x,*}^2:=w^\top H(x)^{-1}w.
\]
The Jeffreys divergence is
\[
    V(x,y):=\ip{x-y}{g(x)-g(y)}.
\]
Let $d_R$ denote the Riemannian distance induced by $H$.

Throughout the numerical examples we use the standard self-concordant normalization
\[
    |\D^3\phi(x)[p,q,q]|
    \le
    2\norm{p}_x\norm{q}_x^2.
\]
Equivalently, the Hessian metric satisfies the exponential comparison
\[
    e^{-2d_R(x,z)}H(x)
    \preceq
    H(z)
    \preceq
    e^{2d_R(x,z)}H(x).
\]
We also use the standard geodesic coercivity inequality
\[
    d_R(x,y)^2\le V(x,y).
\]

For fixed $x,y$, define
\[
    a:=x-y,
    \qquad
    b:=H(x)^{-1}\bigl(g(x)-g(y)\bigr),
\]
so that
\[
    V(x,y)=\ip{a}{b}_x.
\]
Define the asymmetry residual
\[
    U(x,y):=a-b.
\]
This residual vanishes for quadratic/constant-Hessian barriers, and more generally measures the first nonlinear correction to the flat primal-dual matching.

\section{Projected displacement}

Let $L\subseteq \R^n$ be a fixed linear subspace, and let $L^\perp$ denote its Euclidean orthogonal complement. The projected displacement $\eta$ is defined by
\[
    x+\eta\in y+L,
    \qquad
    -g(x)-H(x)\eta\in -g(y)+L^\perp.
\]
Let $P=P_L^x$ be the $H(x)$-orthogonal projection onto $L$, and let $Q=I-P$.

\begin{lemma}[Projection formula]
The projected displacement satisfies
\[
    \eta=-Pb-Qa.
\]
Consequently,
\[
    a+\eta=PU,
    \qquad
    b+\eta=-QU.
\]
\end{lemma}

\begin{proof}
The primal condition $x+\eta\in y+L$ is equivalent to
\[
    a+\eta\in L.
\]
The dual condition says
\[
    -g(x)-H(x)\eta+g(y)\in L^\perp.
\]
Since $g(x)-g(y)=H(x)b$, this is equivalent to
\[
    H(x)(b+\eta)\in L^\perp.
\]
Because $L^\perp$ is Euclidean orthogonal, the last condition is equivalent to
\[
    b+\eta\in L_x^\perp,
\]
where $L_x^\perp$ is the $H(x)$-orthogonal complement of $L$.
Thus $a+\eta$ is the $L$-component of $a-b=U$, while $b+\eta$ is the negative $L_x^\perp$-component of $U$. Hence
\[
    a+\eta=PU,
    \qquad
    b+\eta=-QU.
\]
Equivalently,
\[
    \eta=PU-a=-Pb-Qa.
\]
\end{proof}

The projected energy is
\[
    E:=\norm{\eta}_x^2
    =\norm{Pb}_x^2+\norm{Qa}_x^2.
\]
Moreover, if $x(t)$ is any curve with $x(0)=x$ and $\dot x(0)=\eta$, and $v(t):=V(x(t),y)$, then
\[
    \nabla_R V(x,y)=a+b,
\]
where $\nabla_R$ denotes the Riemannian gradient. Hence
\[
\begin{aligned}
    v'(0)
    &=\ip{a+b}{\eta}_x \\
    &=\ip{a+b}{-Pb-Qa}_x \\
    &=-\ip{a}{Pb}_x-\ip{b}{Pb}_x-\ip{a}{Qa}_x-\ip{b}{Qa}_x.
\end{aligned}
\]
Using the $P/Q$ orthogonality and $\ip{a}{b}_x=V(x,y)$, this simplifies to
\[
    \boxed{v'(0)=-(E+V(x,y)).}
\]
Thus the projected displacement gives the exact descent identity that cancels the flat term in the second-variation analysis.

\section{Original Bregman-balanced integrator analysis}

Let $z_B=x+\Delta$ be the output of the Bregman-balanced integrator with predicted displacement $\eta$. Equivalently, $z_B$ satisfies
\[
    g(z_B)-g(x)=H(x)(2\eta-\Delta).
\]
This is the optimality condition for
\[
    \min_z\left\{
    \frac12\norm{z-(x+2\eta)}_x^2+D_\phi(z,x)
    \right\},
\]
where
\[
    D_\phi(z,x)=\phi(z)-\phi(x)-\ip{g(x)}{z-x}.
\]
The Bregman term keeps the minimizer in the interior of the domain under the usual barrier assumptions.

Set
\[
    \varepsilon:=\Delta-\eta.
\]
Then
\[
    z_B-y=a+\Delta=(a+\eta)+\varepsilon=PU+\varepsilon.
\]
Also,
\[
\begin{aligned}
    H(x)^{-1}\bigl(g(z_B)-g(y)\bigr)
    &=b+H(x)^{-1}\bigl(g(z_B)-g(x)\bigr) \\
    &=b+2\eta-\Delta \\
    &=b+\eta-\varepsilon \\
    &=-QU-\varepsilon.
\end{aligned}
\]
Therefore
\[
\begin{aligned}
    V(z_B,y)
    &=\ip{PU+\varepsilon}{-QU-\varepsilon}_x \\
    &=-\ip{U}{\varepsilon}_x-\norm{\varepsilon}_x^2.
\end{aligned}
\]
Completing the square gives the exact identity
\[
    \boxed{
    V(z_B,y)
    =
    \frac14\norm{U(x,y)}_x^2
    -
    \left\|\varepsilon+\frac12U(x,y)\right\|_x^2.
    }
\]
Consequently,
\[
    \boxed{V(z_B,y)\le \frac14\norm{U(x,y)}_x^2.}
\]
This bound is global for the Bregman-balanced endpoint; no Dikin condition such as $\norm{\Delta}_x<1$ is needed. The reason is that the implicit Bregman equation pairs the primal endpoint correction and the dual endpoint correction with opposite signs, producing the negative square above.

\section{Residual bound and scalar basin for self-concordance}

The exponential Hessian comparison implies the residual estimate
\[
    \boxed{
    \norm{U(x,y)}_x
    \le
    2\bigl(\cosh d_R(x,y)-1\bigr).
    }
\]
Indeed, along a unit-speed geodesic $\gamma:[0,R]\to\operatorname{int}K$ from $x$ to $y$, with $R=d_R(x,y)$,
\[
    U(x,y)
    =
    \int_0^R
    \left[H(x)^{-1}H(\gamma(s))-I\right]\gamma'(s)\,\dd s.
\]
The comparison
\[
    e^{-2s}H(x)\preceq H(\gamma(s))\preceq e^{2s}H(x)
\]
implies, by a spectral calculation,
\[
    \left\|
    \left[H(x)^{-1}H(\gamma(s))-I\right]\gamma'(s)
    \right\|_x
    \le
    e^s-e^{-s}=2\sinh s.
\]
Integrating gives
\[
    \norm{U(x,y)}_x\le \int_0^R 2\sinh s\,\dd s
    =2(\cosh R-1).
\]
Since $d_R(x,y)^2\le V(x,y)$, define
\[
    j(v):=2\bigl(\cosh\sqrt v-1\bigr).
\]
Then
\[
    \norm{U(x,y)}_x\le j(V(x,y)).
\]
Combining this with the Bregman identity yields
\[
    \boxed{
    V(z_B,y)
    \le
    F_B(V(x,y)),
    \qquad
    F_B(v):=\frac14j(v)^2=(\cosh\sqrt v-1)^2.
    }
\]
Numerically,
\[
    F_B(v)\le v
    \quad\text{for}\quad
    0\le v\le 2.6119004634\ldots.
\]
Moreover,
\[
    F_B(v)\le v^2
    \quad\text{for}\quad
    0\le v\le 8.8974963495\ldots.
\]
In particular, on the simple level set $V(x,y)\le 1$,
\[
    \boxed{V(z_B,y)\le V(x,y)^2\le V(x,y).}
\]
Thus the Bregman-balanced integrator has a large certified quadratic basin in the normalized sense.

\section{Exact geodesic as a Bregman-balanced step plus a defect}

Let $\gamma:[0,R]\to\operatorname{int}K$ be the exact geodesic with
\[
    \gamma(0)=x,
    \qquad
    \gamma'(0)=u,
    \qquad
    \norm{\gamma'(s)}_{\gamma(s)}=1,
\]
and let the projected displacement be
\[
    \eta=Ru.
\]
Let
\[
    z_g:=\gamma(R),
    \qquad
    \Delta_g:=z_g-x,
    \qquad
    \varepsilon_g:=\Delta_g-\eta.
\]
For any candidate endpoint $z=x+\Delta$, define its Bregman balance defect by
\[
    \boxed{
    R_{\rm bal}(z)
    :=
    H(x)^{-1}\bigl(g(z)-g(x)\bigr)-\bigl(2\eta-\Delta\bigr).
    }
\]
The Bregman-balanced integrator is exactly characterized by $R_{\rm bal}(z_B)=0$.

For a general endpoint $z=x+\Delta$, with $\varepsilon=\Delta-\eta$, we have
\[
    z-y=PU+\varepsilon,
\]
while
\[
    H(x)^{-1}\bigl(g(z)-g(y)\bigr)
    =-QU-\varepsilon+R_{\rm bal}(z).
\]
Therefore
\[
\begin{aligned}
    V(z,y)
    &=\ip{PU+\varepsilon}{-QU-\varepsilon+R_{\rm bal}(z)}_x \\
    &=
    \frac14\norm{U}_x^2
    -\left\|\varepsilon+\frac12U\right\|_x^2
    +\ip{z-y}{R_{\rm bal}(z)}_x.
\end{aligned}
\]
Thus
\[
    \boxed{
    V(z,y)
    \le
    \frac14\norm{U}_x^2
    +
    \left|\ip{z-y}{R_{\rm bal}(z)}_x\right|.
    }
\]
For the Bregman integrator, the last term vanishes. For the exact geodesic, the last term measures how far the geodesic endpoint is from satisfying the Bregman balance equation.

\section{Bounding the geodesic balance defect}

Let
\[
    H_s:=H(\gamma(s)),
    \qquad
    H_0:=H(x),
    \qquad
    C_s:=\D^3\phi(\gamma(s)).
\]
Define
\[
    W(s):=\gamma'(s)+H_0^{-1}H_s\gamma'(s).
\]
Using
\[
    g(z_g)-g(x)=\int_0^R H_s\gamma'(s)\,\dd s,
    \qquad
    z_g-x=\int_0^R\gamma'(s)\,\dd s,
\]
we get
\[
    R_{\rm bal}(z_g)=\int_0^R\bigl(W(s)-W(0)\bigr)\,\dd s
    =\int_0^R(R-s)W'(s)\,\dd s.
\]
The geodesic equation in Hessian coordinates is
\[
    \gamma''(s)
    =-\frac12H_s^{-1}C_s[\gamma'(s),\gamma'(s),\cdot].
\]
Set
\[
    c_s:=H_s^{-1}C_s[\gamma'(s),\gamma'(s),\cdot].
\]
Then $\gamma''(s)=-c_s/2$. Also,
\[
    \frac{\dd}{\dd s}\bigl(H_s\gamma'(s)\bigr)
    =C_s[\gamma'(s),\gamma'(s),\cdot]+H_s\gamma''(s)
    =\frac12H_sc_s.
\]
Hence
\[
    W'(s)
    =-\frac12c_s+\frac12H_0^{-1}H_sc_s
    =
    \frac12\bigl(H_0^{-1}H_s-I\bigr)c_s.
\]
This is the key geodesic cancellation: the balance defect begins with the factor $H_0^{-1}H_s-I$, which vanishes at $s=0$.

For standard self-concordance, $\norm{c_s}_{\gamma(s)}\le 2$. The exponential Hessian comparison gives
\[
    \left\|\bigl(H_0^{-1}H_s-I\bigr)c_s\right\|_x
    \le
    (e^s-e^{-s})\norm{c_s}_{\gamma(s)}.
\]
Therefore
\[
    \norm{W'(s)}_x
    \le
    2\sinh s.
\]
Consequently,
\[
    \boxed{
    \norm{R_{\rm bal}(z_g)}_x
    \le
    \int_0^R(R-s)2\sinh s\,\dd s
    =2(\sinh R-R).
    }
\]
This is third order in the geodesic length:
\[
    2(\sinh R-R)=\frac13R^3+O(R^5).
\]

We also need a bound on $\norm{z_g-y}_x$. Since
\[
    z_g-y=PU+\varepsilon_g,
\]
we have
\[
    \norm{z_g-y}_x\le \norm{U}_x+\norm{\varepsilon_g}_x.
\]
Moreover
\[
    \varepsilon_g=\int_0^R(\gamma'(s)-\gamma'(0))\,\dd s
    =\int_0^R(R-s)\gamma''(s)\,\dd s.
\]
Using $\norm{\gamma''(s)}_{\gamma(s)}\le 1$ and
\[
    \norm{w}_x\le e^s\norm{w}_{\gamma(s)},
\]
we obtain
\[
    \norm{\varepsilon_g}_x
    \le
    \int_0^R(R-s)e^s\,\dd s
    =e^R-1-R.
\]
Thus
\[
    \boxed{
    \norm{z_g-y}_x
    \le
    \norm{U}_x+e^R-1-R.
    }
\]
Combining the previous estimates gives the exact-geodesic balance-defect bound
\[
    \boxed{
    V(z_g,y)
    \le
    \frac14\norm{U}_x^2
    +
    \bigl(\norm{U}_x+e^R-1-R\bigr)
    2(\sinh R-R).
    }
\]
Here $R=\sqrt E$.

Using the residual bound $\norm{U}_x\le j(v)$ and the sharp energy envelope
\[
    E\le B(v):=
    \left(
    \frac{\sqrt{4v+j(v)^2}+j(v)}{2}
    \right)^2,
    \qquad v:=V(x,y),
\]
one obtains a fully scalar sufficient bound by setting $R=\sqrt{B(v)}$ in the last display.

As $v\downarrow0$, $j(v)=v+O(v^2)$ and $R=O(\sqrt v)$, so the balance-defect term is $O(v^{5/2})$. Therefore
\[
    V(z_g,y)
    \le
    \frac14v^2+O(v^{5/2}).
\]
This explains why exact geodesics and the Bregman-balanced integrator have the same leading square-completion behavior, while the exact geodesic carries an additional higher-order balance defect.

\section{Why the Bregman integrator is well suited to decrease $V$}

The Bregman-balanced integrator is not merely a numerical approximation to the geodesic. It enforces the endpoint relation
\[
    H(x)^{-1}\bigl(g(z_B)-g(x)\bigr)=2\eta-(z_B-x),
\]
which exactly balances the primal and dual endpoint residuals. This gives the identity
\[
    V(z_B,y)
    =
    \frac14\norm{U(x,y)}_x^2
    -
    \left\|\varepsilon+\frac12U(x,y)\right\|_x^2.
\]
Thus the nonlinear endpoint correction appears with a negative square. This is why the Bregman map is naturally suited to decreasing the target divergence $V$.

By contrast, the exact geodesic satisfies the Bregman balance relation only up to the defect $R_{\rm bal}(z_g)$. The geodesic equation makes this defect third order, but it does not make it vanish. Therefore exact geodesic tracing is better for solving the geodesic ODE, while the Bregman-balanced step is better aligned with the discrete decrease of $V$.

This is not a claim that $V(z_B,y)\le V(z_g,y)$ pointwise for every instance. Rather, it is a structural explanation of the certified bounds:
\[
    \text{Bregman step:}
    \qquad
    V(z_B,y)
    \le
    \frac14\norm{U}_x^2,
\]
whereas
\[
    \text{exact geodesic:}
    \qquad
    V(z_g,y)
    \le
    \frac14\norm{U}_x^2
    +
    \text{balance-defect work}.
\]
The additional work term is higher order locally, but can degrade global sufficient basins.

For comparison, the direct exact-geodesic $v''$ estimate gives the scalar bound
\[
    F_{\rm geo}(v)=\frac12B(v)j(v),
\]
with
\[
    F_{\rm geo}(v)\le v
    \quad\text{for}\quad
    v\lesssim 0.7633246011,
\]
and
\[
    F_{\rm geo}(v)\le v^2
    \quad\text{for}\quad
    v\lesssim 0.4190890308.
\]
The Bregman-balanced scalar map
\[
    F_B(v)=(\cosh\sqrt v-1)^2
\]
satisfies
\[
    F_B(v)\le v
    \quad\text{for}\quad
    v\lesssim 2.6119004634,
\]
and, in particular,
\[
    V(x,y)\le1
    \quad\Longrightarrow\quad
    V(z_B,y)\le V(x,y)^2.
\]
These numbers are only sufficient bounds, not optimal basins, but they support the interpretation that the Bregman-balanced endpoint is especially well matched to the Lyapunov quantity $V$.

\end{document}
